\section{ماژول drone\_compressor (DRONE Compressor)}
\label{sec:module_drone_compressor}

\subsection{معرفی و انگیزه}
این مکانیزم بر مبنای مبانی نظری پیشرفته در ادبیات توجه و فشرده‌سازی \cite{eckart1936approximation} توسعه یافته است.

در چارچوب فشرده‌سازی مدل‌های زبانی، مقایسه میان روش‌های آگاه از داده و توزیع فعال‌سازی‌ها (\lr{Data-Aware Compression}) نظیر \lr{AASVD} یا \lr{FWSVD} و روش‌های استاندارد بدون داده مستلزم وجود خطوط مبنای ساختاریافته و هم‌تراز است. بسیاری از مقایسه‌ها به دلیل تفاوت در روش تخصیص رتبه، استراتژی لایه‌بندی یا سازوکارهای مدیریت خطا دچار انحراف تجربی می‌شوند.

ماژول فشرده‌ساز \lr{DRONE Compressor} (برگرفته از ایده تقریب با وزن‌دهی همانی - \lr{Data-aware Low-Rank SVD with Identity Weighting}) به عنوان خط مبنای کنترل‌شده (\lr{Controlled Baseline}) طراحی گردیده است. این ماژول دقیقاً همان معماری لایه‌بندی ساختاری \lr{SVDLinear}، همان قیود تخصیص پویای رتبه و همان آستانه پذیرش خطای بازسازی فروبنیوس مورد استفاده در الگوریتم‌های پیشرفته آگاه از داده را اعمال می‌کند، اما ماتریس وزن‌دهی فعال‌سازی را به ماتریس همانی ($\mathbf{I}$) تقلیل می‌دهد. این رویکرد به ما امکان می‌دهد اثر خالص هم‌ترازی زیرفضاها و وزن‌دهی فعال‌سازی را بدون دخالت سایر متغیرهای پیاده‌سازی به صورت تجربی و دقیق بسنجیم.

\subsection{مبانی نظری و فرمول‌بندی ریاضی}
در روش‌های آگاه از فعال‌سازی، ماتریس وزن با جذر ماتریس کوواریانس فعال‌سازی‌ها مقیاس‌گذاری می‌شود: $\mathbf{W}_{\text{weighted}} = \mathbf{W} \mathbf{C}_X^{1/2}$. در فشرده‌ساز \lr{DRONE}، این فرض مبنایی قرار می‌گیرد که فعال‌سازی‌های ورودی سفید و ناهمبسته بوده و دارای کوواریانس همانی هستند:
\begin{equation}
\mathbf{C}_X = \mathbf{I} \implies \mathbf{W}_{\text{weighted}} = \mathbf{W}
\end{equation}

بدین ترتیب، تجزیه مقادیر منفرد مستقیماً بر روی ماتریس وزن اولیه $\mathbf{W} \in \mathbb{R}^{d_{\text{out}} \times d_{\text{in}}}$ اجرا می‌گردد:
\begin{equation}
\mathbf{W} = \mathbf{U} \mathbf{\Sigma} \mathbf{V}^\top
\end{equation}

رتبه فشرده‌سازی $r$ با استفاده از تابع پویای تخصیص رتبه با اعمال قید کاهش قطعی پارامترها تعیین می‌شود:
\begin{equation}
r = \min\left(r_{\text{target}}, d_{\text{in}}, d_{\text{out}}, \left\lfloor 0.75 \frac{d_{\text{in}} \cdot d_{\text{out}}}{d_{\text{in}} + d_{\text{out}}} \right\rfloor \right)
\end{equation}

مقادیر منفرد $\mathbf{\Sigma}_r$ با جذر متقارن به فاکتورهای چپ و راست تخصیص داده می‌شوند:
\begin{equation}
\mathbf{U}_r = \mathbf{U}_{:, :r} \, \mathbf{\Sigma}_r^{1/2}, \quad \mathbf{V}_r = \mathbf{\Sigma}_r^{1/2} \, \mathbf{V}_{:, :r}^\top
\end{equation}

ماتریس بازسازی‌شده $\mathbf{W}_{\text{recon}} = \mathbf{U}_r \mathbf{V}_r$ مورد ارزیابی خطای فروبنیوس قرار می‌گیرد:
\begin{equation}
\epsilon_{\text{fro}} = \frac{\|\mathbf{W} - \mathbf{W}_{\text{recon}}\|_F}{\max(\|\mathbf{W}\|_F, 10^{-12})}
\end{equation}
در صورتی که $\epsilon_{\text{fro}} > \epsilon_{\text{max}}$ (با آستانه مجاز $\epsilon_{\text{max}} = 0.50$) باشد، فشرده‌سازی لایه لغو شده و ساختار اولیه حفظ می‌گردد. در غیر این صورت، لایه خطی با ماژول دو مرحله‌ای جایگزین می‌شود:
\begin{equation}
\mathbf{y} = (\mathbf{x} \mathbf{V}_r^\top) \mathbf{U}_r^\top + \mathbf{b}
\end{equation}

\subsection{معماری و پیاده‌سازی}
کلاس \lr{DRONECompressor} در فایل \lr{Codes/framework/compression\_engine/drone\_compressor.py} مستقر بوده و تحت شناسه‌های \lr{drone} و \lr{drone\_compressor} در رجیستری فشرده‌سازی ثبت گردیده است. محورهای شاخص پیاده‌سازی عبارتند از:
\begin{itemize}
    \item \textbf{هدف‌گیری اختصاصی پروجکشن‌های توجه (\lr{Target Attention Only}):} به صورت پیش‌فرض، پارامتر \lr{target\_attention\_only} فعال بوده و ماتریس‌های وزن پروجکشن توجه (\lr{q\_proj}، \lr{k\_proj}، \lr{v\_proj}، \lr{o\_proj}، \lr{c\_attn}، \lr{c\_proj}) را هدف قرار داده و از لایه‌های پیش‌خور متوالی صرف‌نظر می‌کند.
    \item \textbf{بهره‌گیری از هسته ماژولار \lr{SVDLinear}:} استفاده مشترک از همان کلاس \lr{SVDLinear} مورد استفاده در \lr{AASVD} برای تضمین برابری شرایط اجرای سخت‌افزاری، تاخیر و مصرف حافظه.
    \item \textbf{پشتیبانی از لایه‌های \lr{Conv1D} در ترنسفورمرهای علّی:} تبدیل هوشمند ابعاد و ترنسپوز صحیح ماتریس‌های وزن در مدل‌هایی نظیر \lr{GPT-2} بدون تخریب ابعاد پنهان.
    \item \textbf{انعطاف در اتصال آداپتورهای باقیمانده \lr{LoRA}:} پیش‌بینی پارامترهای رتبه آداپتور (\lr{lora\_r}) برای مراحل پس‌آموزش و بازیابی دقت در لایه‌های فاکتورگیری‌شده.
\end{itemize}

\subsection{نتایج تجربی}
ارزیابی ماژول \lr{drone\_compressor} بر روی مدل‌های معیار بر اساس فایل‌های خروجی تجربی آزمایشگاهی در \lr{results/} به ثبت رسیده و در جدول \ref{tab:benchmark_drone_compressor} خلاصه شده است:

\begin{table}[H]
\centering
\small
\caption{نتایج بنچ‌مارک سخت‌افزاری و زبانی فشرده‌ساز \lr{drone\_compressor} روی مدل‌های پایه}
\label{tab:benchmark_drone_compressor}
\begin{tabular}{lccccc}
\toprule
\textbf{مدل پایه \lr{(Base Model)}} & \textbf{دقت \lr{(Accuracy)}} & \textbf{شاخص \lr{F1}} & \textbf{حافظه \lr{(Peak VRAM)}} & \textbf{توان \lr{(Throughput)}} & \textbf{وضعیت} \\
\midrule
\lr{DistilBERT} & \lr{90.71\%} & \lr{0.9095} & \lr{368.29 MB} & \lr{99,966.34 tok/s} & \lr{PASS} \\
\lr{RoBERTa} & \lr{94.04\%} & \lr{0.9421} & \lr{514.29 MB} & \lr{14,583.46 tok/s} & \lr{PASS} \\
\lr{GPT-2} & \lr{0.00\%} & \lr{0.0000} & \lr{816.23 MB} & \lr{1,420.86 tok/s} & \lr{PASS\_NO\_COMPRESSION} \\
\lr{Qwen} & \lr{---} & \lr{---} & \lr{---} & \lr{---} & \lr{SKIP\_NOT\_RUN} \\
\bottomrule
\end{tabular}
\end{table}

بررسی نتایج جدول \ref{tab:benchmark_drone_compressor} نکات حائز اهمیتی را نشان می‌دهد:
\begin{itemize}
    \item در مدل‌های انکودر، ماژول \lr{drone\_compressor} با برتری چشمگیر عمل کرده و به دقت خارق‌العاده \lr{94.04\%} و شاخص اف‌وان \lr{0.9421} در \lr{RoBERTa} و دقت \lr{90.71\%} با توان پردازشی \lr{99,966.34 tok/s} در \lr{DistilBERT} دست یافته و نشان قطعی \lr{PASS} را احراز نموده است.
    \item در مدل دیکودر \lr{GPT-2} بر روی ارزیابی زبانی \lr{WikiText-2}، به دلیل فیلتر شدن لایه‌ها بر مبنای شرط \lr{target\_attention\_only} و برابری ابعاد هسته‌ها، هیچ لایه‌ای فشرده نشده و مدل در وضعیت مستند \lr{PASS\_NO\_COMPRESSION} (با پرپلکسیتی ارزیابی‌شده \lr{41.35}) ثبت شده است.
    \item برای مدل مقیاس بزرگ \lr{Qwen}، به دلیل محدودیت‌های حافظه در سرور بنچ‌مارک و اولویت‌بندی آزمون‌های آگاه از فعال‌سازی، آزمایش به عنوان \lr{SKIP\_NOT\_RUN} نشانه‌گذاری شده است.
\end{itemize}

\subsection{محدودیت‌ها و تحلیل شکست}
تحلیل مکانیزم‌های ماژول \lr{drone\_compressor} نشان‌دهنده ابعاد زیر در رفتار الگوریتم است:
\begin{itemize}
    \item \textbf{آسیب‌پذیری در غیاب داده‌های کالیبراسیون برای لایه‌های نامتقارن:} در شرایطی که جهت‌های مقادیر تکین ماتریس وزن هم‌راستا با فعال‌سازی‌های پنهان ورودی نباشد، حذف مؤلفه‌های کوچک تکین بدون لحاظ انرژی فعال‌سازی می‌تواند به از دست رفتن ویژگی‌های حساس زبانی بینجامد؛ هرچند در تسک‌های انکودر متقارن مانند \lr{SST-2} دقت در سطح عالی حفظ شده است.
    \item \textbf{رفتار محافظه‌کارانه در مدل‌های دیکودر علّی (\lr{No Compression}):} الگوریتم به دلیل اعمال قید سخت‌گیرانه کاهش پارامتر نسبت به لایه‌های \lr{Conv1D} و تفکیک سرهای توجه، در مواردی ممکن است از فشرده‌سازی لایه‌ها صرف‌نظر کند که نیازمند تنظیم مجدد پیکربندی لایه‌های هدف است.
    \item \textbf{سربار در فرآیند استنتاج نسبت به حالت متراکم اولیه:} چنانچه رتبه انتخابی به اندازه کافی کوچک نباشد، تفکیک یک لایه خطی متراکم به دو لایه ضرب خطی متوالی ممکن است به دلیل دو بار فراخوانی کرネル محاسباتی و افزایش ترافیک حافظه پنهان، افزایش توان پردازشی متناسبی ایجاد ننماید.
\end{itemize}
